\Chapter{Realizacja przepływu danych}
\label{ch:implementation}

\section{Kontrakty i kontrola granic etapów}

Wspólne walidatory wymagają właściwych typów kolumn, niepustych kluczy
i braku duplikatów. Dla pomiarów dopuszczonych jako nieznane pozostawia
się wartości puste; nieskończoność jest błędem. Numery okrążeń nie mogą
być tekstem lub liczbą zmiennoprzecinkową tylko dlatego, że dałyby się
przekonwertować. Wymuszenie typów po odczycie Parquet ogranicza ryzyko,
że kolejne etapy zinterpretują te same dane inaczej~\cite{contracts}.

Kontrakty obejmują również planowane artefakty. Tabela próbek ma
przechowywać stabilny \code{row_id}, klucz $K$, etykietę, podział i statusy.
Tablice \code{tabular.npz} oraz \code{sequences.npz} muszą mieć te same
identyfikatory wierszy w tej samej kolejności. W ten sposób nawet poprawna
numerycznie macierz zostanie odrzucona, jeśli odpowiada innemu uporządkowaniu
etykiet. Liczba cech wynosi 15, a okno sekwencyjne ma 15 kroków i 16 kanałów:
ostatni kanał \code{is_pad} oznacza dopełnienie z lewej strony.

Dla $N$ próbek interfejs przewiduje $X_{\mathrm{tab}}\in
\mathbb{R}^{N\times15}$ i $X_{\mathrm{seq}}\in
\mathbb{R}^{N\times15\times16}$ w typie \code{float32}.
Walidator sekwencji sprawdza skończoność wartości, binarność i ciągłość
dopełnienia oraz rzeczywisty ostatni krok. Z samej macierzy nie odtworzy
jednak, czy okno zawiera okrążenia jednego kierowcy i kończy się dokładnie
na $t$. Te własności wymagają przyszłych testów pochodzenia okien.

\section{Eksport snapshotu offline}

Funkcja \code{build_snapshot} najpierw sprawdza nazwę wersji, zakres
sezonów i zgodność polityki czyszczenia z D3. Istniejący katalog lub
archiwum nie mogą zostać zastąpione. Obliczenia odbywają się w katalogu
tymczasowym. Kod zapisuje rewizję producenta oraz informację o lokalnych
niezatwierdzonych zmianach, co pozwala odróżnić odtwarzalny kod od
zmodyfikowanego drzewa roboczego.

Kontekst \code{offline_cache} ustawia tryb offline FastF1 oraz blokuje
transport HTTP, transport uwierzytelnienia i połączenia gniazd sieciowych.
Dodatkowo \code{cached_call} obserwuje komunikaty o brakujących
odpowiedziach w cache. Jest to potrzebne, ponieważ biblioteka może
wewnętrznie stłumić błąd ładowania; sam powrót z funkcji nie wystarcza
do stwierdzenia kompletności danych.

Dla każdego wyścigu \code{load_race} pobiera z cache tabelę okrążeń,
status toru, komunikaty kontroli wyścigu i zapowiedzi liczby okrążeń.
\code{raw_laps} normalizuje typy bez usuwania wierszy i sprawdza kontrakt.
\code{coverage} oblicza liczby wierszy, zakresy i wewnętrzne luki okrążeń
oraz braki poszczególnych pól dla obserwowanych kierowców.
Nie dowodzi to obecności kierowcy, którego danych w ogóle nie odczytano;
nie odtwarza także nieobserwowanych początków i końców historii.

Błąd terminarza lub wyścigu pozostaje zapisany w diagnostyce i uniemożliwia
udostępnienie gotowego snapshotu. Po pełnym eksporcie program zapisuje
pliki, wyznacza ich rozmiary i skróty SHA-256 oraz tworzy manifest.
Następnie uruchamia weryfikację, buduje archiwum, rozpakowuje je do nowego
katalogu tymczasowego i sprawdza ponownie. Dopiero po tych kontrolach
przenosi katalog oraz archiwum do docelowych nazw. Osobna blokada plikowa
serializuje tę końcową operację~\cite{project}.

\begin{table}[htbp]
\centering\small
\caption{Zawartość snapshotu T2 i znaczenie czasu.}
\label{tab:snapshot}
\begin{tabularx}{\textwidth}{L{0.33\textwidth}Y}
\toprule
Plik & Zawartość \\
\midrule
\code{laps.parquet} & Surowe okrążenia; czasy w sekundach sesji. \\
\code{events.json} & Metadane wyścigu i ostatnia niepusta liczba planowanych okrążeń. \\
\code{track_status.parquet} & Statusy ze znacznikiem czasu w sekundach sesji. \\
\code{race_control.parquet} & Komunikaty; czas jako tekst ISO w UTC, a nie sekundy sesji. \\
\code{lap_count_evidence.json} & Zapowiedzi w kolejności źródłowej, także aktualizacje puste; materiał archiwalny. \\
\code{manifest.json} & Wersja, producent, pochodzenie, skróty, pokrycie, wykluczenia i błędy. \\
\bottomrule
\end{tabularx}
\end{table}

Ważne jest rozróżnienie podstaw czasu w tabeli~\ref{tab:snapshot}.
Bez przeliczenia czasu UTC komunikatów kontroli wyścigu nie można go
bezpośrednio porównać z sekundami sesji z tabeli okrążeń.
Eksporter zachowuje te dane, ale nie implementuje jeszcze takiego
połączenia na potrzeby cech.

\section{Weryfikacja niezależna od pamięci podręcznej}

\code{verify_data} odczytuje wyłącznie rozpakowany snapshot.
Sprawdza deklaracje plików, rozmiary, SHA-256 i rzeczywistą zawartość
katalogu. Odrzuca ścieżki bezwzględne, próby wyjścia poza katalog wersji
oraz dowiązania symboliczne. Manifest nie znajduje się na własnej liście
skrótów; jego dokładne bajty można identyfikować odrębnym skrótem.

Integralność bajtów nie wystarcza, gdy błędne dane zostały ponownie
zahaszowane. Dlatego weryfikator odczytuje tabele, sprawdza schematy,
tożsamość wyścigów, metadane i znaczniki czasu oraz ponownie oblicza
pokrycie i wykluczenia. Poprawka T2 wymaga także, by każdy zadeklarowany
strumień znajdował się na liście plików. Pola \code{Status} i
\code{Message} muszą zawierać tekst lub braki, a dla każdego wyścigu
co najmniej jedną użyteczną, niepustą wartość. Częściowe braki pozostają
dozwolone. Historyczny układ bez deklaracji opcjonalnych strumieni nie
został automatycznie uznany za niepoprawny.

CLI rozróżnia trzy wyniki: kod 0 oznacza kompletny, poprawny snapshot,
kod 2 wewnętrznie spójne, ale niekompletne pokrycie, a kod 1 błąd formatu
lub integralności. To rozróżnienie pozwala skryptowi wywołującemu
przerwać dalszą pracę bez analizowania swobodnego opisu błędu.

\section{Pochodzenie i pozostały konflikt kontraktowy}

Czas utworzenia snapshotu nie jest czasem pozyskania każdego źródła.
Istniejące cache i dziennik T0 nie uzasadniają odtworzenia dokładnych
czasów pobrań. T2 zapisuje więc \code{acquired_at_utc: null} z wyjaśnieniem,
zamiast nadawać mu czas budowy. Człowiek zaakceptował ograniczony wyjątek
dla istniejących danych, lecz formalna zmiana kontraktu T2c nadal
oczekuje realizacji i poprzedza pełną akceptację snapshotu.

Pozostaje też rozbieżność dokumentacji: D10 dopuszcza FastF1 3.9,
natomiast starszy opis manifestu w chronionym kontrakcie nadal wymienia
3.8.x. Raport opisuje faktyczną wersję i działający weryfikator 3.9,
nie traktując tej rozbieżności jako samodzielnej zmiany decyzji.
Archiwalny plik zapowiedzi liczby okrążeń nie jest zatwierdzonym
interfejsem przyszłego konsumenta D5~\cite{runbook,t2merge}.
